\clearpage
\appendix

\section{Training Procedure}
\label{sec:appendix_algorithm}

Algorithm~\ref{alg:sart} summarizes one training step of SART. The validation gradient $\mathbf{g}_{\mathcal{V}}$ is refreshed every $k$ steps (synthetic: $k{=}5$; LLM: once per epoch) rather than per-step, amortizing the only backprop pass through the validation set.

\begin{figure}[h]
\centering
\small
\begin{tabular}{p{0.90\linewidth}}
\toprule
\textbf{Algorithm~1: SART training step (synthetic scale).}\\
\midrule
\textbf{Input:} batch $\mathcal{B}{=}\{s_i\}$, params $\theta$, val set $\mathcal{V}$, hyperparameters $(\alpha,\beta,\tau_A,\tau_R,\lambda,\gamma,\rho)$, step $t$, refresh interval $k$.\\
1.~If $t \bmod k = 0$: $\mathbf{g}_{\mathcal{V}} \leftarrow \nabla_\theta \tfrac{1}{|\mathcal{V}|}\sum_v \ell(\theta;v)$.\\
2.~For each $s_i \in \mathcal{B}$: compute $\mathbf{g}_{s_i},\ \mathbf{g}_{s_i}^{\mathrm{ans}},\ \mathbf{g}_{s_i}^{\mathrm{reason}}$.\\
3.~Compute $A(s_i){=}\cos(\mathbf{g}_{s_i},\mathbf{g}_{\mathcal{V}})$ and $R(s_i){=}\|\mathbf{g}_{s_i}^{\mathrm{ans}}\|/(\|\mathbf{g}_{s_i}^{\mathrm{ans}}\|{+}\|\mathbf{g}_{s_i}^{\mathrm{reason}}\|)$.\\
4.~Score $S(s_i){=}\alpha\max(0,\tau_A{-}A(s_i))+\beta\max(0,R(s_i){-}\tau_R)$;\ weight $w(s_i){=}\exp(-\lambda S(s_i))$.\\
5.~Aggregate: $\bar{\mathbf{g}}\leftarrow \tfrac{1}{|\mathcal{B}|}\sum_i w(s_i)\,\mathbf{g}_{s_i}$.\\
6.~Conflict-only projection: if $\bar{\mathbf{g}}^\top \mathbf{g}_{\mathcal{V}} < 0$,\ \ $\bar{\mathbf{g}}\leftarrow \bar{\mathbf{g}}-\gamma\,(\bar{\mathbf{g}}^\top\mathbf{g}_{\mathcal{V}}/\|\mathbf{g}_{\mathcal{V}}\|^2)\,\mathbf{g}_{\mathcal{V}}$.\\
7.~Answer-gradient suppression: on samples with $R(s_i){>}\tau_R$, replace $\mathbf{g}_{s_i}^{\mathrm{ans}}$ contribution by $(1{-}\rho)\mathbf{g}_{s_i}^{\mathrm{ans}}{+}\mathbf{g}_{s_i}^{\mathrm{reason}}$ before step~5.\\
8.~$\theta \leftarrow \theta - \eta\,\bar{\mathbf{g}}$.\\
\bottomrule
\end{tabular}
\label{alg:sart}
\end{figure}

\paragraph{Per-sample variant.} For pretrained-LLM scale, Step~6 is promoted to a per-sample gate: for each $s_i$, test $\mathbf{g}_{s_i}^\top \mathbf{g}_{\mathcal{V}} < 0$ and project $\mathbf{g}_{s_i}$ before aggregation. This costs one extra loop pass and removes the averaging that otherwise lets clean gradients outvote shortcut gradients (see Appendix~\ref{sec:appendix_llm_transfer}).

\paragraph{Reasoning-consistency generation protocol.} At evaluation time, we feed the question prefix up to (but excluding) the reasoning-chain separator, greedily generate with temperature $0$ until the answer-token marker, and report the fraction of examples whose generated reasoning tokens exactly match the gold chain up to that marker.

%-----------------------------------------------------------------------------
\section{Graceful Degradation without CoT Supervision}
\label{sec:appendix_graceful}

Eq.~(4) in Section~\ref{sec:method} assumes access to a token-level answer/reasoning partition, which is naturally available for chain-of-thought-supervised datasets (e.g., GSM8K, MATH) and for any dataset where the final answer field can be string-matched back to a contiguous suffix of the target sequence. When such a split is unavailable---for instance, when only the final answer is supervised and no intermediate reasoning tokens exist ($T_{\mathrm{reason}}{=}\emptyset$)---the denominator in Eq.~(4) is no longer informative. In that regime we set $R(s){=}\tau_R$ (equivalently $C(s){=}0$), so that $S(s){=}\alpha\,B(s)$ degenerates to an \emph{alignment-only} ShortcutScore that still delivers the validation-gradient-based sample detection of Step~1. This gives a smooth fallback from a two-signal to a one-signal detector without requiring architectural changes, and lets SART apply to answer-only supervision in the same pipeline. When $T_{\mathrm{reason}}$ can be recovered heuristically (e.g., ``everything before the last numeric span'' for arithmetic tasks), we reinstate Eq.~(4) over the heuristic split.

%-----------------------------------------------------------------------------
\section{Scalable Per-Sample Gradients for Pretrained LLMs}
\label{sec:scalable_gradients}

A naive implementation of SART stores a dense per-sample gradient vector of dimension $d = |\theta|$, which for a 7B-parameter model is $\sim$28\,GB in FP32 and prohibitive to materialize at batch scale. We keep the method tractable on pretrained LLMs through three concrete mechanisms.

\noindent\textbf{LoRA-only gradient subspace.}
We apply SART on top of parameter-efficient fine-tuning: the base model is frozen and trainable capacity is confined to LoRA adapters~\cite{hu2022lora} with rank $r{=}16$ on the attention projections, giving roughly $8{-}16$M trainable parameters for a 7B backbone. Because $\nabla_\theta \ell$ is identically zero on the frozen parameters, the per-sample gradient vector collapses from 28\,GB to $\sim$32\,MB without changing the ShortcutScore or surgery equations in Section~\ref{sec:method}. This also aligns SART with the dominant deployment mode of modern LLMs, where full fine-tuning is rarely feasible.

\noindent\textbf{Random-projection sketches.}
When even the LoRA-gradient footprint is inconvenient---e.g., across hundreds of samples in the scoring pass---we replace each vector by a fixed-seed Gaussian random projection of dimension $d_{\mathrm{sk}}{=}1024$. Under the Johnson--Lindenstrauss lemma, inner products $\langle \mathbf{g}_s, \mathbf{g}_{\mathcal{V}}\rangle$ and hence cosine alignment $A(s)$ are preserved within $\varepsilon$-error with high probability, so the ShortcutScore computed from sketches is a consistent estimator of its dense counterpart. The sketch is applied to $\mathbf{g}_s$, $\mathbf{g}_s^{\mathrm{ans}}$, $\mathbf{g}_s^{\mathrm{reason}}$, and $\mathbf{g}_{\mathcal{V}}$ in the scoring pass; surgery, which operates on dense gradients in-place at the batch level, is unaffected.

\noindent\textbf{Periodic validation-gradient refresh.}
$\mathbf{g}_{\mathcal{V}}$ changes slowly relative to the model parameters, so we recompute it every $k$ optimization steps (synthetic: $k{=}5$; LLM: $k{=}$ once per epoch). With $|\mathcal{V}|{\ll}|\mathcal{D}|$, this amortizes the cost of the only step that backpropagates through the validation set.

These three mechanisms are orthogonal to the ShortcutScore and surgery formulations: any alternative proxy that preserves inner-product geometry---Fisher-diagonal reweighting, last-layer gradients, or implicit-differentiation approximations~\cite{koh2020understandingblackboxpredictionsinfluence}---can be slotted in without altering Eqs.~(5) or~(\ref{eq:gradient_surgery}).

%-----------------------------------------------------------------------------
\section{Hyperparameter Sensitivity under Conflict-Only Surgery}
\label{sec:appendix_sensitivity}

\begin{table}[h]
\centering
\caption{Hyperparameter sensitivity for SART's full method under the conflict-only (PCGrad-gated) surgery rule. Results are per-dataset because the optimal projection strength $\gamma$ depends on the compositional complexity of the shortcut; Avg reports the uniform mean across the three synthetic benchmarks.}
\label{tab:sensitivity}
\small
\begin{tabular}{llcccc}
\toprule
\textbf{$(\gamma, \rho)$} & \textbf{Metric} & \textbf{Math} & \textbf{Financial} & \textbf{Causal} & \textbf{Avg} \\
\midrule
\multirow{2}{*}{$(0.8, 0.7)$ — prior default}
    & Acc ($\uparrow$) & \textbf{85.4} & 67.6 & 55.1 & 69.4 \\
    & Rob ($\uparrow$) & \textbf{47.1} & 43.9 & 28.3 & 39.8 \\
\midrule
\multirow{2}{*}{$(0.8, 0.3)$}
    & Acc ($\uparrow$) & \textbf{85.4} & 67.6 & --- & --- \\
    & Rob ($\uparrow$) & \textbf{47.1} & 43.9 & --- & --- \\
\midrule
\multirow{2}{*}{$(1.0, 0.3)$ — recommended}
    & Acc ($\uparrow$) & 84.2 & \textbf{74.8} & \textbf{68.6} & \textbf{75.9} \\
    & Rob ($\uparrow$) & 22.5 & \textbf{65.7} & \textbf{63.3} & \textbf{50.5} \\
\midrule
\multirow{2}{*}{per-task best $(\gamma, \rho)$}
    & Acc ($\uparrow$) & \textbf{85.4} & \textbf{74.8} & \textbf{68.6} & \textbf{76.3} \\
    & Rob ($\uparrow$) & \textbf{47.1} & \textbf{65.7} & \textbf{63.3} & \textbf{58.7} \\
\bottomrule
\end{tabular}
\end{table}

Table~\ref{tab:sensitivity} reports SART's full method (reweighting + conflict-only surgery + answer suppression) at three $(\gamma, \rho)$ settings on each synthetic benchmark, plus the per-task best (the dataset-wise argmax over the table). The sensitivity structure that emerges is materially different from what the prior (non-gated) surgery rule suggested.

\textbf{(1) $\gamma$ should scale with the compositional complexity of the shortcut.} Math-Arithmetic has a single-feature shortcut (first operand $\geq 5$); here $\gamma{=}0.8$ is best, and pushing $\gamma$ to $1.0$ over-projects and collapses robustness from 47.1 to 22.5. Financial-Analysis and Causal-Reasoning both have multi-feature true rules (conjunctions and confounder adjustment); here $\gamma{=}1.0$ is best, with $+7.2$\,pp accuracy and $+21.8$\,pp robustness on Financial, and $+13.5$\,pp accuracy and $+35.0$\,pp robustness on Causal. The underlying driver is that single-feature shortcuts leave informative residual signal along $\mathbf{g}_{\mathcal{V}}$ that full projection would destroy, while multi-feature shortcuts require full removal of the conflicting subspace.

\textbf{(2) $\rho$ is inactive unless $R(s)$ crosses $\tau_R$.} On Math and Financial the answer-gradient concentration $R(s)$ rarely exceeds $\tau_R{=}0.5$, so $\rho$ in $[0.3, 0.7]$ yields identical numbers (confirmed by the $(0.8, 0.3)$ row). Only on Causal, whose compositional answer tokens produce high-concentration gradients, does $\rho$ matter: a Causal mini-sweep shows $\rho{=}0.3$ and $\rho{=}0.0$ are tied (both at 68.6 / 63.3), while $\rho{=}0.5$ drops to 55.1 / 28.3. The previous recommendation of $\rho{=}0.5$ was a pre-fix artifact from counter-balancing the unconditional projection rule.

\textbf{(3) The recommended default $(\gamma{=}1.0, \rho{=}0.3)$ improves over the prior default on average.} Against the earlier $(\gamma{=}0.8, \rho{=}0.7)$ setting, the recommended default gains $+6.5$\,pp accuracy and $+10.7$\,pp robustness on average; per-task tuning adds another $+8.2$\,pp robustness by switching to $\gamma{=}0.8$ on Math. In deployment, a practitioner who cannot tune per-task should use $(\gamma{=}1.0, \rho{=}0.3)$; one who can should estimate the shortcut's feature cardinality (e.g., from the shortcut detector's coefficient structure) and back off $\gamma$ for single-feature cases.

%-----------------------------------------------------------------------------
\section{Why Detection F1 Decouples from Downstream Gains}
\label{sec:appendix_f1_decoupling}

Table~\ref{tab:ablation} shows the Reweighting-Only variant reaching the highest standalone detection F1 (0.85), yet the full SART pipeline---which acts on a noisier F1 of 0.64---delivers substantially larger accuracy and robustness gains. We trace this decoupling to the difference between \emph{how a detection signal is acted on} rather than \emph{how accurate it is}.

Reweighting-Only computes sharp per-sample weights with a single decision threshold, so the F1 cleanly reflects the classifier's separability, but the correction is confined to gradient \emph{magnitude}. The full method uses a smoother exponential weight together with conflict-only projection: the score sweep that computes F1 now mixes along-$\mathbf{g}_{\mathcal{V}}$ and across-$\mathbf{g}_{\mathcal{V}}$ effects in a single scalar, lowering the oracle F1 even though neither false-positive (weight slightly reduced on a benign sample) nor false-negative (uncorrected conflict still survives projection when $\mathbf{g}_s^\top \mathbf{g}_{\mathcal{V}}\ge 0$) is catastrophic. Data Filtering exhibits the opposite imbalance: its F1 of 0.69 is achieved by \emph{dropping} the detected samples, which in high-FP regimes discards genuine reasoning signal wholesale---a cost that is invisible to the F1 metric but shows up directly in downstream accuracy. Higher detection F1 is therefore necessary but not sufficient; what matters is whether the action conditioned on the detector (drop vs.\ soft-reweight vs.\ surgery) preserves the information the detector got right while bounding the damage from the information it got wrong.

%-----------------------------------------------------------------------------
\section{Real-World LLM Transfer on GSM8K}
\label{sec:appendix_llm_transfer}

To test whether SART's gradient-level shortcut correction transfers beyond the synthetic benchmarks, we conduct a controlled-injection pilot on GSM8K with Qwen3-8B using LoRA rank-16 adapters on $q/k/v/o\text{-proj}$ layers (15.3M trainable parameters, $0.19\%$ of the backbone). Shortcut is the sum of all numbers in the question; 70\% of training samples are relabeled to the shortcut answer and the model is evaluated on the original true-answer test split. The setup is designed so that a capable base model would otherwise fully internalize the shortcut: standard fine-tuning collapses to $0.5\%$ clean accuracy with the generation pattern ``Adding the values: $x{+}y{+}z = s$. \texttt{\#\#\#\#} $s$'' on essentially every test item, confirming there is meaningful headroom for a correction mechanism to claim.

\begin{table}[h]
\centering
\caption{Ablation on Qwen3-8B / GSM8K (70\% shortcut injection). Clean accuracy and clean perplexity on 200 held-out items. PPL is the primary indicator of whether the intervention preserves generation; accuracy indicates whether it breaks the shortcut.}
\label{tab:llm_transfer}
\small
\begin{tabular}{lccccl}
\toprule
\textbf{Variant} & $\gamma$ & $\rho$ & \textbf{Acc.\ (\%)} & \textbf{PPL} & \textbf{Outcome} \\
\midrule
SFT (baseline)                          & --  & --  & 0.5 & 1.56                 & memorizes shortcut \\
SART, default                           & 0.5 & 0.3 & 1.0 & $2.76{\times}10^{5}$ & generation collapsed \\
SART, $\rho{=}0$                        & 0.5 & 0.0 & 1.0 & 1.54                 & generation ok, shortcut persists \\
SART, reweight-only                     & 0   & 0   & 0.5 & 1.55                 & surgery inert at batch level \\
SART, $\gamma{=}1$                      & 1.0 & 0.0 & 0.5 & 1.54                 & $\gamma$ is batch-level no-op \\
SART, per-sample PCGrad                 & 1.0 & 0.0 & 0.5 & 1.54                 & gate fires, shortcut persists \\
SART, per-sample + $\lambda{=}5$        & 1.0 & 0.0 & 0.5 & 1.53                 & heavy reweighting insufficient \\
\bottomrule
\end{tabular}
\end{table}

\textbf{Detection transfers; correction does not, across the full hyperparameter cross-section tested.} Phase~2 ShortcutScore separates the 70\% injected samples (mean $S = 0.46$) from the 30\% clean samples (mean $S = 0.007$) by more than an order of magnitude---the same kind of separation seen on the synthetic benchmarks. The correction step, however, does not move accuracy under any configuration in Table~\ref{tab:llm_transfer}. We trace this to two mechanical issues that surface when SART is applied unmodified to an autoregressive LLM:

\emph{(1) Answer-token suppression is incompatible with autoregressive generation.} With $\rho > 0$ the suppression term subtracts gradient from the exact tokens that teach the model to emit the final numeric answer, and perplexity explodes by five orders of magnitude ($1.56 \to 2.76\times 10^{5}$). Setting $\rho = 0$ restores generation quality (PPL $\leq 1.55$) but of course disables the suppression component entirely.

\emph{(2) Batch-level conflict gating averages away the shortcut signal.} Conflict-only projection (Eq.~\ref{eq:gradient_surgery}) is implemented on the batch-aggregated gradient $\mathbf{g}_{\mathrm{batch}}$: projection fires when $\mathbf{g}_{\mathrm{batch}}^\top \mathbf{g}_{\mathcal{V}} < 0$. On mixed batches containing both shortcut and clean samples, the clean samples' reasoning gradients align positively with $\mathbf{g}_{\mathcal{V}}$ and outvote the negatively-aligned shortcut gradients, so the gate almost never fires. Doubling the projection strength from $\gamma=0.5$ to $\gamma=1.0$ changes nothing because the gate condition is independent of $\gamma$; the two configurations produce identical per-sample training loss curves and identical accuracy/PPL on held-out data.

Promoting the gate from batch-level to per-sample---tested in the fifth row---is the minimal fix that preserves SART's original theoretical argument (project only when the sample itself conflicts with the validation gradient, PCGrad-style) while removing the mixing artifact. Implementation-wise, it adds a single inner loop over the mini-batch at $\sim\!2{\times}$ the cost of the batch-level path. The per-sample gate does fire: the Phase~3 training loss stays $\sim\!20\%$ higher than under batch-level surgery at the same $\gamma$, indicating that gradient components are actually being projected. The test-time accuracy, however, still tracks SFT. The last row tests whether stronger reweighting closes the gap by raising $\lambda$ from 1 to 5 (mean shortcut sample weight drops from $\approx 0.64$ to $\approx 0.11$, and 97\% of training samples are now heavily down-weighted); it does not. Across the full hyperparameter cross-section tested, correction does not move accuracy.

We read this as a cleanly negative correction result at 8B scale under a 3-epoch LoRA budget: \emph{given} that SFT fully collapses to the shortcut on this dataset, the gradient-level interventions SART specifies---soft reweighting and conflict-gated projection---are mechanically sound (all instrumentation reports the expected behavior) but insufficient to overturn the optimization asymmetry between a compact shortcut (``sum all numbers'') and multi-step arithmetic on a capable base model. Three concrete follow-ups fall out of the mechanism analysis above: (i) restrict $\mathbf{g}_s^{\mathrm{ans}}$ to the \emph{numeric} answer tokens so that $\rho > 0$ can be re-enabled without killing generation; (ii) promote detection to a hard cut (drop samples with $S(s) > \tau$) so that 100\% of the remaining training signal is shortcut-free rather than the $\lambda=5$ floor of $\approx 11\%$ contamination; (iii) expand the trainable parameter budget (LoRA rank or additional layers) and the training budget (10--20 epochs at low LR) so that the non-shortcut solution has room to out-fit the shortcut. The detection-side transfer is a clean positive result that we keep separate from any subsequent correction-side claim.
